% Capítulo incluído por tcc-ufv.tex com \include: sem preâmbulo e sem \begin{document}.
% Na classe abntex2-ufv, \autoref de um capítulo imprime "Seção N" (o trabalho é dividido em seções).
\chapter{Introdução}\label{cap:introducao}

A elaboração de trabalhos acadêmicos segue as normas da \ac{abnt}. Este
modelo mostra as peças de um trabalho completo com a classe
\texttt{abntex2-ufv}, cada capítulo em um arquivo.

\section{Objetivos}\label{sec:objetivos}

O objetivo geral é desenvolver e avaliar um protótipo. Os objetivos
específicos são:
\begin{alineas}
  \item levantar os requisitos (\autoref{cap:metodologia});
  \item desenvolver o protótipo;
  \item avaliar o protótipo com usuários (\autoref{cap:resultados}).
\end{alineas}

\section{Organização do trabalho}

A \autoref{cap:referencial} apresenta os conceitos usados no trabalho; a
\autoref{cap:metodologia}, os métodos; a \autoref{cap:resultados}, os
resultados; e a \autoref{cap:conclusao}, as conclusões.
